\chapter{Evaluation}
\label{chap:Evaluation}

% \TBD is defined in implementation.tex; redefined here only if this chapter is
% compiled on its own.
\providecommand{\TBD}[1][]{\textcolor{red}{[TBD\if\relax\detokenize{#1}\relax\else: #1\fi]}}
\providecommand{\placeholderfig}[1]{%
  \fbox{\parbox[c][4.5cm][c]{0.85\linewidth}{\centering\small\color{red} #1}}}

This chapter reports six experiments. The first motivates the black-box assumption.
The next three test hypotheses H1 to H3 of Section~\ref{sec:Problem Statement} in
order, since each is worth testing only if the previous one holds. The last two
analyse the selection rewards and the allocations that the selectors learn. For each
hypothesis the decision rule is stated before the results. Before these, a pilot run with
one seed, reported in Section~\ref{eval:sec:Pilot}, tested the pipeline and identified the
revisions needed for the main runs.

\section{Hardware}
\label{sec:Hardware}

The main experiments run on a shared server with six NVIDIA L40 GPUs of 48\,GB each,
two AMD EPYC 7763 processors and 1\,TB of memory, without a job scheduler. Jobs wait for
an idle GPU, use at most two GPUs at a time and one GPU per run. Providers and the
student's rollouts are served with vLLM, and training uses PyTorch with the PEFT library
(Section~\ref{imp:sec:config}). Screening the candidate providers took 3.0 GPU-hours and
building the response cache 8.1 GPU-hours, and the pilot run of
Section~\ref{eval:sec:Pilot} 14.4 GPU-hours. The motivating experiment of
Section~\ref{eval:sec:Motivating} runs on the School of Computing's Slurm cluster, on
NVIDIA A100 and H100 GPUs. \TBD[total GPU-hours once the experiments are complete]

\section{Experimental Setup}
\label{sec:Experimental Setup}

\paragraph{Metrics.} The primary metric is test utility $U(\theta_\tau)$, the
weighted mean of per-domain test accuracy, with equal domain weights unless stated.
Every result is reported together with the budget spent and the number of gradient
steps taken, since at equal budget cheaper providers buy more steps
(Section~\ref{imp:sec:config}). Where students of different capacity are compared, the
teacher-gap recovery
\[
  R = \frac{\mathrm{Acc}_{\mathrm{method}} - \mathrm{Acc}_{\mathrm{base}}}
           {\mathrm{Acc}_{\mathrm{teacher}} - \mathrm{Acc}_{\mathrm{base}}}
\]
normalises away the student's starting point.

\paragraph{Budgets.} Each selector is run at three budgets $B \in \{\TBD, \TBD, \TBD\}$,
chosen so that the smallest allows roughly \TBD\ rounds with the strongest provider
alone.

\paragraph{Seeds and uncertainty.} Every configuration is run with three seeds (0, 1
and 2). Intervals are 95\% bootstrap intervals over seeds and test items. Selectors
are compared by paired differences on the same seeds and test items.

\section{Pilot Run}
\label{eval:sec:Pilot}

Before the main experiments, every condition was run once, with seed 0 and the single
budget $B = 2{,}400$, to test the pipeline end to end and to check the design against
real training. The pilot comprised the five single-provider runs, uniform allocation,
the headroom selector, discounted Thompson sampling, the provider-free correctness
control and the SFT baseline, together with an evaluation of the untrained student. All
eleven jobs completed, using 14.4 GPU-hours. Each final student was evaluated greedily
on the 500-problem test set of every domain.

With one seed, these results describe what happened rather than test the hypotheses.
Each mean in Table~\ref{eval:tab:pilot} is an accuracy over 1,500 test problems, with a
standard error of about 0.013, so differences in mean below about 0.03 are not
distinguishable. The decision rules of Sections~\ref{eval:sec:H1} to~\ref{eval:sec:H3}
are not applied here.

\begin{table}[htbp]
\centering
\small
\begin{tabular}{@{}lrrcccc@{}}
\toprule
 & & & \multicolumn{4}{c}{\textbf{Test accuracy}} \\
\cmidrule(l){4-7}
\textbf{Condition} & \textbf{Rounds} & \textbf{Spent} & \textbf{Countdown} & \textbf{MATH} &
  \textbf{Code} & \textbf{Mean} \\
\midrule
Untrained student            & --  & --    & 0.218 & 0.434 & 0.336 & 0.329 \\
\midrule
Qwen2.5-3B only              & 400 & 1,380 & 0.682 & 0.426 & 0.368 & 0.492 \\
Phi-4-mini only              & 400 & 2,057 & 0.714 & 0.440 & 0.350 & 0.501 \\
Qwen2.5-7B only              & 353 & 2,397 & 0.694 & 0.426 & 0.352 & 0.491 \\
Qwen2.5-Coder-7B only        & 394 & 2,392 & 0.304 & 0.456 & 0.350 & 0.370 \\
Phi-4 only                   & 109 & 2,377 & 0.508 & 0.424 & 0.336 & 0.423 \\
\midrule
Uniform                      & 217 & 2,397 & 0.604 & 0.434 & 0.338 & 0.459 \\
Headroom per cost            & 400 & 1,798 & 0.584 & 0.476 & 0.338 & 0.466 \\
Discounted Thompson sampling & 185 & 2,399 & 0.504 & 0.428 & 0.358 & 0.430 \\
\midrule
Correctness reward, no provider & 200 & -- & 0.560 & 0.440 & 0.376 & 0.459 \\
SFT on provider responses    & --  & 2,384 & 0.378 & 0.432 & 0.328 & 0.379 \\
\bottomrule
\end{tabular}
\caption[Pilot results]{Pilot results with one seed. Rounds are GRPO steps; runs that
reached the cap of 400 rounds stopped before spending the whole budget. Spent is in
price units; the pre-query filter leaves skipped prompts uncharged.}
\label{eval:tab:pilot}
\end{table}

\paragraph{Learning happened almost only on Countdown.} Every on-policy run raised
Countdown accuracy, from 0.218 to as much as 0.714. On MATH and code, no condition moved
test accuracy by more than 0.042 from the untrained student, which is within about 1.3
standard errors. The training logs show why. During training the student's sampled
accuracy on pool prompts did rise, from about 0.37 to 0.43--0.48 on MATH and from about
0.23 to 0.32--0.36 on code in the four single-provider runs that spent more than 100
rounds on each domain, but only to roughly the level of its existing greedy accuracy
(0.47 on MATH and 0.40 on code on the selection set). The
updates made the student's sampling more consistent with what it already answered
greedily, rather than teaching it to solve new problems. On Countdown, by contrast,
sampled accuracy rose from about 0.1 to 0.6--0.7 under Qwen2.5-3B, Phi-4-mini and
Qwen2.5-7B, and greedy accuracy followed. At this
scale of training, the selection signal on MATH and code is therefore close to zero for
every arm, and H1 and H2 can only be observed on Countdown. A follow-up check trained
on each domain alone for 60 rounds with Qwen2.5-7B, at the pilot's learning rate and at
five times that rate. The higher rate did not help on MATH ($-0.014$ against $+0.022$ on
the selection set) or code ($+0.000$ against $+0.010$), cut the gain on Countdown from
$+0.280$ to $+0.098$, and shortened the student's Countdown responses from 361 to 169
tokens. The learning rate therefore stays at $2 \times 10^{-5}$, and MATH and code are kept
as domains in which the student learns little at this scale; a good selector should learn
not to spend there.

\paragraph{Providers differ in how fast they teach Countdown.}
Figure~\ref{eval:fig:pilot-countdown} follows Countdown accuracy in the single-provider
runs. Per round, Phi-4-mini taught fastest, reaching 0.71 after 100 Countdown rounds,
against 0.60 for Qwen2.5-7B and 0.51 for Qwen2.5-3B. Per unit of budget the ordering
changed as spending grew: Qwen2.5-3B led up to about 200 price units, and Phi-4-mini was
ahead by 400 (0.62 against 0.53). This is the shape that H1 and H2 describe, but here it
is confounded by the rotation between domains and rests on one seed; the branch
measurements of Section~\ref{eval:sec:H2} are the test. Phi-4 could afford only about 40
Countdown rounds within the budget. The provider-free control, rewarded with ground-truth
correctness, reached 0.54 on the selection subsample after 60 Countdown rounds, against 0.55 for Phi-4-mini and 0.46
for Qwen2.5-3B, so on this task agreement with a provider was about as useful as the
true labels.

\begin{figure}[htbp]
\centering
\includegraphics[width=\linewidth]{pilot_countdown_learning.pdf}
\caption[Countdown accuracy in the single-provider pilot runs]{Countdown accuracy in the
single-provider pilot runs, against the number of Countdown rounds (left) and the budget
spent on Countdown (right). Accuracy is measured greedily on the 128-problem selection
subsample after each block, so single points carry a standard error of about 0.04.}
\label{eval:fig:pilot-countdown}
\end{figure}

\paragraph{One provider made the student collapse.} Under Qwen2.5-Coder-7B, the
student's Countdown responses shrank from about 235 tokens of reasoning to about 20
tokens, a bare answer, within about 40 Countdown rounds, and its Countdown accuracy on
the selection subsample stayed below 0.3. No other run showed this: in every other run
the student's Countdown responses ended between 445 and 584 tokens long. The cause is not
yet clear. Invalid provider answers alone do not explain it, since 64\% of
Qwen2.5-Coder-7B's cached Countdown responses contain no valid answer, against a similar
66\% for Phi-4-mini, which caused no collapse. A provider can therefore harm the student,
not merely help it less. A second explanation was also tested and rejected: a bare guess
that adds the three numbers would be rewarded if the provider often answered that way, but
Qwen2.5-Coder-7B's most common answer is that trivial sum on only 2.9\% of prompts, no more
often than for the other providers (2 to 8\%). The higher learning rate of the follow-up
check also shortened the student's responses, which suggests that the collapse is a failure
mode of strong updates rather than a property of one provider. The main runs test whether
it recurs with other seeds.

\paragraph{The headroom selector made the predicted switch, then stalled on MATH.}
Figure~\ref{eval:fig:pilot-headroom} shows its run. It began with Qwen2.5-3B on Countdown,
as Table~\ref{imp:tab:predictions} predicts. Countdown accuracy reached 0.297 after round
40, just past the predicted switch point of 0.27, and from round 41 the selector moved to
Qwen2.5-7B. It then turned to Qwen2.5-3B on MATH and spent 320 of its 400 rounds there.
MATH accuracy did not rise, so the headroom on MATH never shrank and the selector rarely
left. Countdown, the domain that was learning, received only 60 rounds. When MATH
accuracy briefly reached 0.48 around round 120, its headroom fell far enough for
Qwen2.5-7B on Countdown to win twice, and those two blocks raised Countdown accuracy from
0.40 to 0.52. Its two blocks with the degraded provider (rounds 11--20 and 161--170)
were random exploration: at Phi-4's price and with a lower ceiling, the degraded provider
never has the highest value. The run reached the round cap having spent 1,798 of its 2,400 units.
The assumption that progress is proportional to headroom
(Section~\ref{imp:sec:headroom}) fails when a domain has headroom but no learnability, and
the selector has no feedback from realised gains to correct it. The main runs use a
revised selector that learns this feedback (Section~\ref{imp:sec:headroom-learned}).

\begin{figure}[htbp]
\centering
\includegraphics[width=\linewidth]{pilot_headroom_timeline.pdf}
\caption[The headroom selector in the pilot]{The headroom selector in the pilot.
Top: greedy accuracy on the selection subsample per domain. Bottom: the arm chosen in
each block of ten rounds, coloured by provider.}
\label{eval:fig:pilot-headroom}
\end{figure}

\paragraph{Thompson sampling chased noise and favoured expensive arms.} It spent 1,326
of its 2,399 units, 55\%, on the degraded provider. Two causes are visible in its log.
First, it had only 20 blocks in which to learn about 18 arms. Second, every block's
measured gain lay between $-0.023$ and $+0.029$, which is within the noise of comparing
two evaluations on 128 problems per domain (about $\pm 0.035$). Dividing such noisy gains
by cost magnifies the noise of cheap arms and shrinks that of expensive ones towards
zero: a block of Qwen2.5-3B on MATH costing 34 units scored $-0.69$ per thousand units,
while blocks of the degraded provider costing about 210 units scored between $-0.07$
and $+0.14$. With
gains this small, cost normalisation biases the selector towards expensive arms whatever
their value. The main runs use a cost-aware version that estimates gain per round and
divides by price only when choosing (Section~\ref{imp:sec:selectors}), and every block is
now evaluated on 256 problems per domain instead of 128.

\paragraph{Uniform allocation spent most of its budget on the expensive providers.}
Because it gives every arm the same number of blocks, 1,582 of its 2,397 units (66\%) went
to Phi-4 and its degraded copy, which explains why it trailed the cheap single-provider
runs despite visiting every arm.

\paragraph{Off-policy fine-tuning helped least.} SFT on provider responses, bought with
the same budget, reached 0.378 on Countdown against 0.604 for the on-policy uniform
allocation, and did not improve MATH or code.

\paragraph{The pre-query filter saved 15.7\% of queries.} Across all runs, 15.7\% of the
prompts drawn for training were skipped before the provider was consulted, from 6.7\% in
the Phi-4 run to 23.5\% in the Phi-4-mini run. It was highest in the runs where the
student's answers had converged most, on Countdown under Phi-4-mini and in the collapsed
Qwen2.5-Coder-7B run.

\paragraph{What the pilot implies for the hypotheses.} With one seed, no adaptive
selector beat the best single provider: the headroom selector reached a mean of 0.466 and
Thompson sampling 0.430, against 0.501 for Phi-4-mini alone. This is not a test of H3,
because both selectors showed design flaws that are fixable before the main runs. The
pilot does show that the best provider for a domain can differ from the most capable
one, and that the most cost-effective provider can change as the student learns, which is
the premise of H1 and H2.

\paragraph{Revisions for the main runs.} The main runs keep the learning rate, replace
both selectors with the revised versions of Sections~\ref{imp:sec:headroom-learned}
and~\ref{imp:sec:selectors}, and evaluate each block on 256 problems per domain. In a
simulator calibrated to the pilot, the revised selectors no longer showed the pilot's
failures (Section~\ref{imp:sec:verification}). The single-provider runs, uniform
allocation, provider-free control and SFT baseline are unchanged and are rerun with three
seeds.

\section{Motivating Experiment: The Cost of Cross-Tokenizer Supervision}
\label{eval:sec:Motivating}

\paragraph{Question.} How much of the teacher--student gap does white-box on-policy
distillation recover when teacher and student do not share a tokenizer, and why?
The answer determines whether the black-box assumption gives anything up.

\paragraph{Setup.} The teacher is Qwen3-4B-Instruct. Students are Qwen2.5-1.5B-Instruct,
which shares the teacher's tokenizer, OLMo-2-1B, with vocabulary overlap 0.65, and
SmolLM2-1.7B, with overlap 0.24. All runs start from each student's SFT checkpoint
on verified teacher completions and train on Countdown with the same trainer and
step count; only the loss and the student change.

\begin{table}[H]
\centering
\begin{tabular}{|c|l|l|c|c|}
\hline
\textbf{ID} & \textbf{Student} & \textbf{Method} & \textbf{Test acc.} & $R$ \\
\hline
B1 & each & SFT on teacher completions & \TBD & 0 \\
T  & teacher & -- & \TBD & 1 \\
G1 & Qwen2.5-1.5B & GKD, exact logits & \TBD & \TBD \\
G2 & Qwen2.5-1.5B & ULD only, same tokenizer & \TBD & \TBD \\
X1 & OLMo-2-1B & GOLD hybrid & \TBD & \TBD \\
X2 & SmolLM2-1.7B & GOLD hybrid & \TBD & \TBD \\
\hline
\end{tabular}
\caption[Same-tokenizer and cross-tokenizer distillation on Countdown]{Countdown test
accuracy and teacher-gap recovery for same-tokenizer and cross-tokenizer on-policy
distillation.}
\label{eval:tab:motivating}
\end{table}

\paragraph{Result.} \TBD[report $R$ for G1 versus X1 and X2, the G1--G2 gap as the cost
of the loss approximation alone, and the diagnostics that support each failure mode]

\section{Experiment 1: Non-Dominance (H1)}
\label{eval:sec:H1}

\begin{mdframed}[backgroundcolor=bg, linewidth=0pt, innertopmargin=6pt,
  innerbottommargin=6pt]
\textbf{Decision rule.} H1 is supported if, at $\theta^{(0)}$, there are two domains
$d \ne d'$ whose arms of highest estimated marginal value belong to different
providers, and in each the leading arm's 95\% interval lies above that of the other
provider's arm.
\end{mdframed}

\paragraph{Protocol.} Marginal values $\hat v(k, d; \theta^{(0)})$ are estimated by
branching from the initial student, as in Section~\ref{imp:sec:Marginal}.

\begin{table}[H]
\centering
\begin{tabular}{|l|c|c|c|}
\hline
\textbf{Provider} & \textbf{Countdown} & \textbf{MATH} & \textbf{Code} \\
\hline
Phi-4                   & \TBD & \TBD & \TBD \\
Qwen2.5-7B              & \TBD & \TBD & \TBD \\
Qwen2.5-Coder-7B        & \TBD & \TBD & \TBD \\
Qwen2.5-3B              & \TBD & \TBD & \TBD \\
Phi-4-mini              & \TBD & \TBD & \TBD \\
Phi-4, degraded         & \TBD & \TBD & \TBD \\
\hline
\end{tabular}
\caption{Estimated marginal value $\hat v(k, d; \theta^{(0)})$, in utility points per
unit cost, with 95\% intervals. Bold marks the best provider in each domain.}
\label{eval:tab:h1}
\end{table}

A second table reports the same quantities without dividing by cost, to separate a
provider that is most useful from one that is most useful per unit of budget.

\paragraph{Result.} \TBD

\section{Experiment 2: State Dependence (H2)}
\label{eval:sec:H2}

\begin{mdframed}[backgroundcolor=bg, linewidth=0pt, innertopmargin=6pt,
  innerbottommargin=6pt]
\textbf{Decision rule.} H2 is supported in domain $d$ if the provider of highest
estimated marginal value at some checkpoint $\theta^{(c)}$ differs from that at
$\theta^{(0)}$, with non-overlapping 95\% intervals between the two providers at
$\theta^{(c)}$. H2 is supported overall if it holds in at least one domain under
the on-policy update.
\end{mdframed}

\paragraph{Protocol.} The branching of Section~\ref{imp:sec:Marginal} is repeated at
each checkpoint $\theta^{(1)}, \dots, \theta^{(C)}$. Rank stability is summarised by
Kendall's $\tau$ between the provider rankings at consecutive checkpoints in each
domain.

\begin{figure}[H]
\centering
\placeholderfig{Placeholder: one panel per domain; $x$-axis checkpoint (budget
spent); $y$-axis $\hat v(k, d; \theta^{(c)})$ with intervals; one line per provider.
A crossing of lines is a change of the dominant provider.}
\caption{Marginal value of each provider across training, per domain.}
\label{eval:fig:h2}
\end{figure}

\paragraph{Result.} \TBD

\paragraph{If H2 fails.} A negative result shows that a single allocation fixed
before training suffices on these tasks. Experiment~3 is then reported only to
confirm that adaptive selectors do no better than the best static allocation.

\section{Experiment 3: Does Adaptivity Pay? (H3)}
\label{eval:sec:H3}

\begin{mdframed}[backgroundcolor=bg, linewidth=0pt, innertopmargin=6pt,
  innerbottommargin=6pt]
\textbf{Decision rule.} H3 is supported if, at some budget, the best deployable
adaptive selector exceeds the best static allocation in hindsight, with the 95\%
interval of the paired difference above zero.
\end{mdframed}

\begin{table}[H]
\centering
\begin{tabular}{|l|c|c|c|}
\hline
\textbf{Selector} & $B_{\mathrm{low}}$ & $B_{\mathrm{mid}}$ & $B_{\mathrm{high}}$ \\
\hline
Uniform                    & \TBD & \TBD & \TBD \\
Strongest-only             & \TBD & \TBD & \TBD \\
Cheapest-only              & \TBD & \TBD & \TBD \\
Best static in hindsight   & \TBD & \TBD & \TBD \\
\hline
Headroom, learned efficiency (proposed) & \TBD & \TBD & \TBD \\
Cost-aware Thompson sampling & \TBD & \TBD & \TBD \\
Cost-aware EXP3            & \TBD & \TBD & \TBD \\
Sliding-window UCB         & \TBD & \TBD & \TBD \\
Contextual (LinUCB)        & \TBD & \TBD & \TBD \\
\hline
Greedy oracle              & \TBD & \TBD & \TBD \\
\hline
\end{tabular}
\caption{Test utility at each budget under the on-policy update, mean and 95\%
interval over seeds. A second table reports the same under the SFT update.}
\label{eval:tab:h3}
\end{table}

\begin{figure}[H]
\centering
\placeholderfig{Placeholder: test utility (y) against budget spent (x), one line per
selector, shaded intervals.}
\caption{Utility as a function of budget spent.}
\label{eval:fig:h3}
\end{figure}

\paragraph{Result.} \TBD

\section{Experiment 4: Selection Rewards}
\label{eval:sec:Rewards}

\paragraph{Question.} Can the cheap proxies of Table~\ref{imp:tab:rewards} stand in for
utility gain, which requires an
evaluation pass after every block of rounds?

\paragraph{Protocol.} On the branches of Experiment~2, the rank correlation between
each proxy and the measured utility gain is computed per domain. The best adaptive
selector from Experiment~3 is then rerun with each proxy in place of utility gain.

\begin{table}[H]
\centering
\begin{tabular}{|l|c|c|c|}
\hline
\textbf{Selection reward} & \textbf{Spearman $\rho$ with utility gain}
  & \textbf{Final utility} & \textbf{Extra compute} \\
\hline
Utility gain   & 1 & \TBD & \TBD \\
Loss decrease  & \TBD & \TBD & \TBD \\
Learnability   & \TBD & \TBD & \TBD \\
Signal density & \TBD & \TBD & \TBD \\
\hline
\end{tabular}
\caption{Agreement of each selection reward with utility gain, and the final utility
a selector reaches when driven by it.}
\label{eval:tab:rewards}
\end{table}

\paragraph{Result.} \TBD

\section{Experiment 5: Allocation Analysis and Robustness}
\label{eval:sec:Allocation}

\paragraph{Question.} How does each selector divide its budget among the providers?
Does its spending track each provider's contribution, and does it learn to avoid the
degraded provider?

\paragraph{Protocol.} For each selector, the share of budget spent on each provider is
compared with a post hoc valuation of that provider's contribution to final utility
\autocite{sim_cml_2020}, computed by \TBD[leave-one-provider-out retraining, or a
sampled Shapley estimate over provider subsets]. The fraction of budget spent on the
degraded provider is tracked over time, and the degraded provider's corruption rate
$p$ is varied over \TBD.

\begin{figure}[H]
\centering
\placeholderfig{Placeholder: left, stacked area of budget share per provider over
rounds for the best adaptive selector; right, budget share against post hoc value,
one point per provider, with the identity line.}
\caption{Budget allocation over training and its relation to provider value.}
\label{eval:fig:allocation}
\end{figure}

\paragraph{Result.} \TBD

\section{Summary of Findings}
\label{eval:sec:Summary}

\TBD[one paragraph per hypothesis: supported or not, by how much, and under which
update rule]

